\documentclass{article}

% Issue #117: only one half of an alias pair is defined, because the other half
% needs no definition — the literal `\begin{X}`/`\end{X}` is a spelling of each
% side. `\bsplit` expands to `\begin{split}`, so a written-out `\end{split}`
% closes it; `\eeq` expands to `\end{equation}`, so it closes a written-out
% `\begin{equation}`. This is the shape the reporter's paper is full of.
\def\bsplit{\begin{split}}
\def\eeq{\end{equation}}

\begin{document}

% `split` is math-only, so the shape that actually ships is nested inside a
% display environment — the alias dispatch has to reach math mode to see it.
\begin{equation}
  \bsplit
    a &= b, \\
    c &= d.
  \end{split}
\end{equation}

% The mirror, on the other side of the same equation.
\begin{equation}
  a = b
\eeq

% Both spellings live at once, each pairing with the shape it is.
\begin{equation}
  \bsplit
    e &= f
  \end{split}
\eeq

% Gated shapes: none of these pair, and none of them diagnose. Ordered so each
% one's would-be closer is genuinely out of reach — the last `\end{split}` in
% the file is the one on the math line, so the opener below it has none at all.

% A closer stranded outside the brace group the opener sits in.
\foo{\bsplit stranded}
\end{split}

% A closer behind a math delimiter the scan does not model.
\bsplit $x$ behind math \end{split}

% No reachable closer whatsoever.
\bsplit

\end{document}
